%%==========================================================
%%  CAPÍTULO 5: DISEÑO Y CONSTRUCCIÓN DEL GESTOR DE MÓDULOS REUTILIZABLES
%%==========================================================
\label{cap:gestor}

El Capítulo~\ref{cap:generacion} dejó una herramienta que produce un proyecto
conforme al marco y un contrato por el que una capacidad compartida puede
acoplársele. Lo que no existía todavía era esa capacidad, ni la operación que
la incorpora. Un proyecto generado nace completo el primer día y empieza a
quedarse corto el segundo, cuando necesita algo que no eligió al derivarse, y
esta fase construyó lo que resuelve ese segundo día.

Lo que entregó son dos cosas que se validan mutuamente, y la primera son
módulos reales del catálogo, que es la única forma de comprobar que el
contrato de integración sirve para algo. La segunda son las operaciones que
los incorporan, los consultan, los actualizan y los retiran sobre un proyecto
que ya existe y que su equipo ya modificó. La condición de llegada se formuló
antes de construir, y exige que un producto derivado al que se le incorpora un
módulo siga arrancando sin edición manual, supere su propia verificación,
conserve el código que su equipo escribió y declare qué módulo tiene y en qué
versión. El detalle de la construcción se conserva en los documentos de la
Fase 3, disponibles en la carpeta \texttt{Implementación/} del repositorio del
trabajo de grado.

\section{Hitos de la construcción}
\label{sec:hitos-fase-3}

La fase se ejecutó en cuatro hitos, y el orden entre ellos importa por la
misma razón que en la fase anterior, porque un módulo real contra la
definición encuentra lo que ninguna lectura de la definición encuentra. Por
eso el activo componible y su prueba contra un caso real preceden a las
operaciones que lo manejan, y la evidencia del circuito completo cierra la
fase. La Tabla~\ref{tab:hitos-gestor} recoge qué dejó cada hito y cómo quedó
comprobado.

\begin{table}[H]
\centering
{\footnotesize
\begin{tabular}{p{3.1cm}p{5.9cm}p{5.4cm}}
\toprule
\textbf{Hito} & \textbf{Qué dejó construido} & \textbf{Cómo quedó comprobado} \\
\midrule
Qué es un módulo & El manifiesto con que un módulo se describe a sí mismo, la frontera entre lo empaquetado y lo generado, y las condiciones que lo hacen válido & La herramienta incorpora un módulo sin conocer su nombre, leyendo solo lo que el manifiesto declara \\
El primer módulo & La persistencia relacional, con su puerto en el contrato y su adaptador en el producto & El producto compuesto guarda un dato, se reinicia y lo vuelve a leer \\
Las operaciones & Incorporar, consultar, actualizar y retirar, sobre el punto de extensión de la fase anterior & El núcleo creció sin que ninguna operación anterior cambiara \\
El segundo módulo & La observabilidad, de otra naturaleza que la primera & Entró sin cambiar el núcleo, el esquema del manifiesto ni la operación \\
\bottomrule
\end{tabular}
}
\caption{Hitos de la construcción del gestor de módulos}
\label{tab:hitos-gestor}
\end{table}

La frontera con la fase siguiente es la que separa construir de medir. Esta
fase comprueba que el mecanismo funciona; si funciona en manos de un equipo
que no lo construyó, y cuánto esfuerzo ahorra, es lo que mide el
Capítulo~\ref{cap:evaluacion}.

\section{Las capacidades transversales y la que se construyó}
\label{sec:capacidades-transversales}

El inventario de capacidades partió del análisis cualitativo del contexto y
no de una taxonomía externa. Aquí hay una desviación respecto del
anteproyecto que debe declararse, y es que aquel documento pedía categorizar
las capacidades transversales tomando como referencia los dominios funcionales
del estándar de negocio, y no se hizo, porque ese estándar modela capacidades
bancarias y no infraestructura técnica de aplicaciones. Categorizar un
registro de eventos contra un dominio de negocio habría producido una
clasificación sin poder explicativo.

De las capacidades inventariadas se eligió una, con criterios escritos antes
de aplicarlos: que fuera verificable por comportamiento observable, que no
dependiera de infraestructura externa al entorno de verificación, que
ejercitara el contrato de composición en su forma completa, que su ausencia
fuera visible en el producto, y que no exigiera decisiones de seguridad que
el contexto no hubiera fijado.

La elegida fue la persistencia relacional del producto, y la
elección corrigió una afirmación anterior del propio proyecto, que daba por
decidido que el módulo comprometido sería el de cifrado y manejo de tokens.
El criterio descarta el cifrado precisamente por el último punto, ya que
exige fijar antes decisiones de seguridad, como la gestión de claves, su
rotación y su custodia, que pertenecen a la organización y no a una línea de
producto.

El inventario fijó además un orden de crecimiento, y la primera capacidad de
ese orden se construyó después como segundo módulo
(sección~\ref{sec:segundo-modulo}). Esa segunda construcción es la que
permite afirmar algo que con un solo módulo no era afirmable.

\section{Qué es un módulo}
\label{sec:anatomia-modulo}

Un módulo de esta línea es un paquete publicado que implementa un puerto
técnico del \emph{contract package} y que se describe a sí mismo. La
herramienta no conoce ningún módulo por su nombre, sino que lee un manifiesto
que viaja dentro del paquete y que declara qué puerto provee, qué
configuración necesita, qué archivos aporta al producto y en qué puntos de
este hay que insertar su cableado.

La frontera entre lo que viaja en el paquete y lo que se escribe dentro del
producto tiene una consecuencia que vale más que la elegancia, porque
incorporar un módulo nuevo al catálogo no exige publicar una versión nueva de
la herramienta. Sin ella, cada módulo habría requerido código propio en el
generador, y el catálogo habría quedado acoplado a su versión.

\subsection{La frontera entre lo empaquetado y lo generado}

La distinción que gobierna todo el mecanismo es cuál de las dos partes de un
módulo vive dónde. La \emph{implementación} viaja empaquetada y versionada en
el módulo, y es código que el equipo del producto no edita y que una
actualización reemplaza por completo. El \emph{cableado} se genera dentro del
producto, y comprende el adaptador que ata el puerto a la implementación, la
configuración que el producto valida junto con la suya, y el registro de qué
quedó incorporado.

De ahí se sigue la regla más restrictiva de la fase: un módulo aporta
plantillas, nunca clases base. Si el módulo exportara una clase de
la que el producto heredara, el código del producto quedaría atado a la forma
interna del módulo y una actualización podría romperlo sin tocarlo. Lo
generado es del equipo que lo recibe, que es lo que RNF-02 exige.

Para que el cableado pueda entrar, y sobre todo para que pueda reemplazarse
después, el producto trae puntos de inserción declarados: unos son marcas de
una sola línea donde se agrega algo, y otros son regiones delimitadas cuyo
contenido se sustituye entero. La diferencia importa, porque lo que se
inserta sin delimitar no se puede reemplazar, y eso apareció al construir la
actualización y no al diseñar la composición.

Un módulo escribe también en la documentación del producto, en una región que
él declara. La razón es que el producto dice de sí mismo cosas que dejan de
ser ciertas al componerlo. Un producto que advierte que pierde sus datos al
detenerse deja de perderlos cuando entra el módulo de persistencia, y esa
frase tiene que corregirse por construcción y no por memoria de quien compone.

\subsection{A qué forma de producto sirve}

El manifiesto declara además a qué forma de producto sirve el módulo, y
mientras la línea tuvo una sola plantilla, esa declaración habría sido
redundante; cuando el almacén admitió una segunda forma dejó de serlo, porque
el cableado de un módulo aterriza en puntos de inserción concretos y una
forma que no los tenga lo recibiría roto.

La comparación se hace contra el registro del producto y nunca contra la
plantilla que lo derivó, porque esa plantilla puede no estar instalada, y se
hace por identificador y versión entera y nunca por rango, ya que una versión
de forma se mueve exactamente cuando un punto se mueve o desaparece. Ser
indulgente ahí convertiría un rechazo que nombra la causa en un producto roto
sin aviso.

Las dos declaraciones que deciden una incorporación caben en cinco líneas del
manifiesto, y son el puerto que el módulo provee y la forma a la que sirve.
Todo lo que la herramienta necesita saber de un módulo para aceptarlo o
rechazarlo está ahí, y nada de ello está en la herramienta.

El mensaje del rechazo nombra las dos formas, la que el módulo sirve y la que
el producto tiene, y explica que la incompatibilidad no es cuestión de
esfuerzo sino de que el cableado no tendría dónde aterrizar. El producto queda
sin tocar.

\section{El módulo de persistencia como caso de referencia}
\label{sec:modulo-persistencia}

El módulo se construyó sobre el motor embebido que la plataforma de ejecución
ya trae, sin dependencias de ejecución nuevas. Esa restricción no es
estética, porque una dependencia de ejecución en un entorno regulado abre una
auditoría, y el objetivo era probar el mecanismo de composición y no la
elección de un motor.

La persistencia es además el caso que obliga a ejercitar a la vez las dos
clases de puerto que el Capítulo~\ref{cap:generacion} definió. El repositorio
del agregado es un puerto de dominio, que el producto declara y realiza. El
acceso al motor es un puerto técnico, que el \emph{contract package} declara y
el módulo realiza. El adaptador del repositorio, que vive en el producto,
consume el puerto técnico del módulo, y ninguna dependencia une directamente
al producto con el módulo.

El módulo aporta la capacidad técnica pero no una base de adaptador que el
producto extienda, porque la herencia entre un activo del catálogo y el
código del producto crearía un acoplamiento que ninguna de las dos partes
controla. La consecuencia es que el dominio no conoce el motor de
persistencia, y cambiarlo no altera ninguna declaración de la capa de
dominio. De aquí se desprende una regla adicional de validez, según la cual
un producto no puede tener puertos de dominio sin adaptador, y es la que
después impide retirar un módulo que dejaría al producto incompleto.

\paragraph*{Un hallazgo del contenedor de inyección}

Un módulo no alcanzaba el puerto de configuración del producto, porque el
contenedor de dependencias resuelve dentro del módulo que pide. La solución
fue que el producto exponga los puertos de la línea en un módulo global
propio, separado de su raíz de composición. Es una consecuencia del marco que
ninguna lectura del marco habría anticipado, porque apareció al ejecutar, y es
el tipo de hallazgo que justifica construir un módulo real en lugar de
especificar cómo sería.

\section{Las cuatro operaciones del gestor}
\label{sec:gestor-modulos}

Las operaciones del gestor se montan sobre el punto de extensión que la fase
anterior dejó, y ninguna lo modifica. El núcleo sí creció, y la distinción
entre crecer y modificarse es la que mantuvo la frontera entre las dos fases.

\subsection{Incorporar un módulo a un proyecto existente}
\label{sec:composicion}

La composición resuelve el módulo contra el catálogo, lo obtiene del
registro, lee su manifiesto y convierte lo que el manifiesto declara en un
plan de cambios sobre un producto que ya existe. El plan sigue siendo el
mismo dato inspeccionable de la fase anterior; lo que cambió es que ahora
puede escribir sobre un destino que no está vacío.

Escribir sobre un proyecto ajeno exigió tres cosas que generar no
necesitaba: un modo que admite un destino existente, dos clases de cambio
—reescribir un archivo completo y reemplazar el contenido de una región
delimitada— y respaldo de todo lo que se toca antes de tocarlo. El diseño
original del plan de cambios ya había anticipado la necesidad, lo que
permitió extenderlo en lugar de reemplazarlo.

\paragraph*{Lo que la operación se niega a hacer}

Antes de producir un solo cambio, la operación comprueba nueve reglas sobre
el producto y sobre el módulo. Rechaza cuando el módulo no sirve a la forma
del producto, cuando el destino no es un producto derivado, cuando el activo
no está en el catálogo o no está activo, cuando el módulo ya está
incorporado, cuando otro módulo ya provee el mismo puerto, cuando las
versiones de contrato no coinciden, cuando falta un punto de inserción que el
cableado necesita, y cuando un archivo que el módulo genera ya existe en el
producto.

Que las nueve se comprueben antes y no durante la escritura es lo que hace
que un rechazo deje el producto exactamente como estaba, y esa propiedad se
verifica. Un rechazo que dejara el producto a medio modificar sería peor que
no tener la operación. La razón de rechazar en lugar de intentar arreglar es
que cada regla describe un producto que la herramienta no sabe componer sin
suponer algo sobre código que no escribió, y suponer en código ajeno es el
modo más caro de equivocarse, porque el daño aparece después, en un producto
que ya nadie relaciona con el comando que lo tocó.

\paragraph*{Lo que apareció al ejecutar}

Los tres defectos de la Tabla~\ref{tab:defectos-composicion} no aparecieron
al diseñar la composición sino al correrla, y son la parte más instructiva de
la fase.

\begin{table}[H]
\centering
{\footnotesize
\begin{tabular}{p{4.9cm}p{9.8cm}}
\toprule
\textbf{Defecto} & \textbf{Qué lo causaba} \\
\midrule
Un archivo tocado dos veces perdía el primer cambio & El plan aplicaba cada cambio sobre el archivo original y no sobre el resultado del anterior \\
Un marcador se confundía con su propia mención en la documentación del código & La búsqueda del punto de inserción no distinguía la marca de un comentario que la nombra \\
El producto compuesto compilaba pero no pasaba su análisis estático & Una importación quedaba sin uso al reemplazar una región \\
\bottomrule
\end{tabular}
}
\caption{Defectos detectados al ejecutar la composición}
\label{tab:defectos-composicion}
\end{table}

El tercero es el que mejor justifica el circuito de verificación completo,
porque un producto compuesto que compila pero no pasa su análisis estático es un
producto que su equipo no puede entregar, y ninguna prueba unitaria de la
herramienta lo habría mostrado.

\subsection{Consultar, actualizar y retirar}
\label{sec:consulta-actualizacion}

La consulta responde qué ofrece el catálogo, qué tiene incorporado el
producto, en qué versión, cuál quedó por detrás de la disponible y qué archivo
generado cambió el equipo. Es lo que permite decidir si actualizar antes de
intentarlo, y dice además lo que el producto \emph{tiene} y no lo que podría
tener, que es una corrección posterior. Antes declaraba disponible un módulo
que la composición habría rechazado, porque no comprobaba la forma del
producto.

La actualización mueve un módulo incorporado a una versión posterior sin
intervención manual, y su problema real es qué hacer con lo que el equipo
cambió. La respuesta fue comparar una huella del archivo generado contra la
que el registro guardó: si coinciden, el archivo se reemplaza; si no, se
preserva y la herramienta lo dice. No se fusiona nada, porque fusionar
exigiría que la herramienta decidiera sobre código que no escribió, y la
decisión de informar en lugar de resolver es lo que mantiene el control en el
equipo. La actualización comprueba además tres de las nueve reglas de la
composición, que al principio no repetía, de modo que un cambio de versión
podía mover un módulo a una versión construida contra otra mayor del contrato,
o contra otra forma de producto, sin que nadie lo mirara.

A las tres operaciones se sumó después una cuarta, y es la única pieza del
artefacto que ningún requerimiento pide. Se construyó por decisión
de alcance del autor, declarada como tal, con el argumento de que una
decisión que no se puede deshacer no es una decisión, porque un equipo que no
puede retirar un módulo no lo incorpora para probarlo. La operación lee
únicamente el registro del producto, restaura la región que el módulo había
reemplazado y conserva lo que el equipo editó; y rechaza antes de tocar nada
si al retirarlo un puerto de dominio quedaría sin adaptador. Queda un límite
declarado, porque un archivo generado que el equipo modificó se conserva por
la misma razón que lo conserva la actualización, y ese archivo sigue
importando el módulo que acaba de salir. La herramienta lo dice en lugar de
corregirlo.

\section{El segundo módulo del catálogo}
\label{sec:segundo-modulo}

Un módulo construido demuestra que el mecanismo funciona para ese módulo, pero
no que sea general, porque con un solo caso no hay forma de distinguir un
contrato de composición de una solución escrita para él. El segundo módulo es
la evidencia de esa diferencia, y por eso se eligió de otra naturaleza que el
primero.

La capacidad construida fue la observabilidad, que era la primera
del orden de crecimiento declarado en el inventario. Donde el primer módulo
persiste datos, este los \emph{exporta}, entregando a un colector externo las
entradas del registro, el tramo de traza de cada petición y su duración, en
un protocolo abierto de telemetría y sin dependencias de ejecución nuevas.
Quedó fuera de su alcance exponer un endpoint para que un sistema externo
recolecte métricas, y la razón es del marco y no del módulo, porque un módulo
aporta adaptadores de salida y un endpoint es superficie de entrada.

\paragraph*{Qué quedó probado}

Un módulo de otra naturaleza entró sin cambiar el núcleo de la herramienta, ni
el esquema del manifiesto, ni la operación de composición. Lo único que creció
fue el catálogo, con una entrada, que es la afirmación que el primer módulo no
podía sostener.

El contrato ganó además un puerto con una propiedad que ninguno anterior
tenía, porque es el único que ningún adaptador de la plantilla
satisface. El producto reporta sus señales exactamente igual tenga o no el
módulo incorporado, y lo inyecta de forma opcional, de modo que incorporar el
módulo no cambia una línea del producto. La plantilla tampoco reemplaza su
registro de eventos cuando el módulo está presente, porque el asiento es del
producto, sino que entrega al módulo la misma entrada que ya escribió.

\paragraph*{Dos hallazgos que solo aparecen al ejecutar}

El primero es una regla del marco que nadie había enunciado, y es que un
módulo no puede recibir el puerto al que la plantilla le reporta. El adaptador
pedía el puerto de registro de eventos para escribir sus propias trazas, lo
que cerró un ciclo con el registro del producto; el resultado fue que el
producto arrancó, atendió peticiones y \emph{no exportó nada, en silencio},
que es la peor forma de fallar. La corrección dejó una prueba de regresión en
la plantilla, porque una regla cuya violación no se comprueba no está
verificada sino escrita.

El segundo corrigió una afirmación del propio marco, que se sostenía que una
versión menor del contrato llega a los consumidores sin intervención, y en el
tramo de versiones anterior a la primera estable eso es falso, porque el
rango que un paquete declara no acepta la menor siguiente. La consecuencia
fue que el primer módulo tuvo que republicarse sin un solo cambio de código.
Es exactamente el costo que el marco había anticipado al versionar el
contrato como una unidad, apareciendo por primera vez en ejecución.

\section{Verificación y conformidad}
\label{sec:verificacion-composicion}

La verificación de esta fase no podía ser una suite de pruebas unitarias. Lo
que había que comprobar es un circuito que atraviesa varias herramientas y
dos procesos: derivar un producto, incorporarle el módulo, verificar que el
producto sigue pasando su propia verificación, arrancarlo, almacenar un dato,
detener el proceso, volver a arrancarlo y leer el dato de vuelta. Eso se
construyó como un circuito automatizado que corre en la puerta de
verificación de la línea.

El circuito comprueba además cada rechazo, dejando el producto idéntico en
todos ellos, la actualización preservando el archivo que el equipo modificó,
y el producto volviendo a ser conforme cuando se le retira el último módulo.
Incorpora el mismo módulo a dos productos distintos, con agregados distintos,
y comprueba que el cableado se nombra por el agregado de cada uno, que es la
evidencia de que un activo se comparte entre proyectos sin acoplarlos. Y
porque el segundo producto se compone después de actualizar el primero, los
dos quedan con versiones distintas del mismo módulo a la vez, que es el
criterio literal de RNF-06.

La evidencia se cita en los documentos de la Fase 3 con su fecha y con el
comando que la reproduce; este capítulo no la reproduce, porque una
transcripción de salidas en la tesis queda desactualizada en la primera
versión publicada.

El criterio de la fase anterior se aplicó a todo lo que esta agregó, y dice
que una regla cuya violación no se haya comprobado no está verificada. Cada regla
nueva se rompió a propósito, se comprobó que el circuito fallara con su
propio mensaje, y se restauró.

\section{Resumen del capítulo}

Esta fase cerró el artefacto, de modo que al terminar existe una herramienta
que genera un proyecto conforme al marco, le incorpora una capacidad
transversal sobre un proyecto ya en curso, informa de su estado, la actualiza
preservando lo que su equipo escribió y la retira si hace falta.

Tres decisiones explican el resultado, y la primera es que el módulo se
describe a sí mismo, lo que desacopla el catálogo de la versión de la
herramienta y hace que incorporar un módulo nuevo no sea un cambio de código.
La segunda es la frontera entre la implementación empaquetada y el cableado
generado, con la regla de que un módulo aporta plantillas y nunca clases
base, que es lo que hace que lo generado sea del equipo y que una
actualización no pueda romper código que no reemplaza. La tercera es que la
actualización informa en lugar de fusionar, lo que mantiene en el equipo toda
decisión sobre código que la herramienta no escribió.

La fase dejó además una lección metodológica que ninguna lectura del marco
habría producido, porque varias de sus decisiones corrigen defectos que
aparecieron al \emph{ejecutar} y no al diseñar. El más incómodo y el más útil
lo dio el segundo módulo, donde un adaptador pedía el puerto de registro de
eventos al que la propia plantilla le reporta, cerró un ciclo, y el producto
arrancó,
atendió peticiones y no exportó nada en silencio. Un fallo silencioso es el
que ninguna lectura del marco produce y ninguna prueba del módulo aislado
detecta, y solo apareció al componer el módulo en un producto real y mirar
qué salía por el otro lado.

Ese segundo módulo es también lo que convierte una afirmación en evidencia, porque
con un módulo construido, el mecanismo funciona para ese módulo; con uno de
otra naturaleza, que exporta en lugar de persistir y que entró sin cambiar el
núcleo, ni el esquema del manifiesto, ni la operación de composición, lo que
queda probado es el mecanismo. Es la misma clase de evidencia que la segunda
plantilla aporta del lado de la derivación, y las dos juntas son lo que
permite hablar de una línea de productos y no de un generador con un ejemplo.

Lo que queda no es construcción sino medición, es decir cuánto del esfuerzo de
configuración inicial se reduce realmente, si la composición funciona en
manos de un equipo que no construyó la herramienta, y si un proyecto derivado
permite empezar a escribir lógica de negocio antes que el proceso actual. Es
el objeto del Capítulo~\ref{cap:evaluacion}.
